\bigskip
\bigskip
\subsubsection*{Խնդիրներ և հարցեր թեմա 17-ի վերաբերյալ}

\begin{enumerate}[label=\thesection.\arabic*.]

% 17.1
\item Դիցուք $(X, \tau_1)$ և $(X, \tau_2)$ տոպոլոգիական տարածություններն այնպիսին են, որ $\tau_1 \subset \tau_2$: Ճի՞շտ է պնդումը.
\begin{enumerate}
    \item[ա)] եթե $(X, \tau_1)$-ը կոմպակտ է, ապա $(X, \tau_2)$-ը կոմպակտ է;
    \item[բ)] եթե $(X, \tau_2)$-ը կոմպակտ է, ապա $(X, \tau_1)$-ը կոմպակտ է:
\end{enumerate}

% 17.2
\item Կոմպակտ է արդյոք աջից կիսաբաց ինտերվալների $(\mathbb{R}, \to)$ տարածությունը:

% 17.3
\item Կոմպակտ է արդյոք ա) $[1, 2)$; բ) $[1, 2]$ ենթաբազմությունը $(\mathbb{R}, \to)$ տարածությունում:

% 17.4
\item Գոյություն ունի՞ $(\mathbb{R}, \text{սովոր.})$ տարածության
\begin{enumerate}
    \item[ա)] $(-1, 1)$ ինտերվալի;
    \item[բ)] $[-1, 1]$ հատվածի
\end{enumerate}
անընդհատ արտապատկերում $\mathbb{R}$-ի վրա:

% 17.5
\item Ապացուցեք.
\begin{enumerate}
    \item[ա)] էվկլիդեսյան $\mathbb{R}^{n+1}$ տարածության $S^n$ միավոր սֆերան և $B^{n+1}$ միավոր փակ գունդը կոմպակտ տարածություններ են;
    \item[բ)] $\mathbb{R}P^n, n \ge 0$ պրոյեկտիվ տարածությունները կոմպակտ տարածություններ են:
\end{enumerate}

% 17.6
\item Ապացուցեք. Տոպոլոգիական տարածության ամեն մի երկու կոմպակտ ենթաբազմությունների միավորումը կոմպակտ ենթաբազմություն է:

% 17.7
\item Ի տարբերություն նախորդ խնդրի պնդման, երկու կոմպակտ ենթաբազմությունների հատումը կարող է չլինել կոմպակտ ենթաբազմություն:
\par Բերենք օրինակ. Դիտարկենք իրական թվերի $A$ ենթաբազմությունը կազմված բոլոր ամբողջ թվերից և $\pm \frac{1}{2}$ թվերից՝ $A = \mathbb{Z} \cup \{\pm \frac{1}{2}\}$: Սահմանենք տոպոլոգիա $A$ բազմության վրա համարելով բաց ենթաբազմություններ $\emptyset$-ը և $A$-ն, $\mathbb{Z}$-ի բոլոր ենթաբազմությունները և $A$-ի այն բոլոր ենթաբազմությունները, որոնց լրացումները վերջավոր են: Ցույց տվեք.
\begin{enumerate}
    \item տոպոլոգիայի 1-3 աքսիոմները բավարարվում են;
    \item $A$-ի $B = \mathbb{Z} \cup \{-\frac{1}{2}\}$ և $C = \mathbb{Z} \cup \{\frac{1}{2}\}$ ենթաբազմությունները կոմպակտ են;
    \item $B \cap C$ ենթաբազմությունը կոմպակտ չէ:
\end{enumerate}

% 17.8
\item Ճի՞շտ է պնդումը. Ցանկացած մետրիկական տարածության ամեն մի փակ և սահմանափակ ենթաբազմություն կոմպակտ ենթաբազմություն է:

\begin{hint}
Որպես հակաօրինակ դիտարկեք որևէ $X$ անվերջ բազմություն $\rho(x, y) = \{ 0, \text{երբ } x=y; 1, \text{երբ } x \ne y \}$ մետրիկայով և նրա $B(x_0, 1) = \{x; \rho(x_0, x) \le 1 \}$ ենթաբազմությունը, որտեղ $x_0$-ն $X$-ի որևէ սևեռված կետ է:
\end{hint}

% 17.9
\item Ապացուցեք. Եթե $A$-ն $X$ հաուսդորֆյան տարածության կոմպակտ ենթաբազմություն է, ապա $A$-ի և ցանկացած $x_0 \in X \setminus A$ կետի համար գոյություն ունեն չհատվող շրջակայքեր:

\begin{hint}
Օգտվել թեորեմ 5-ի ապացուցման ընթացքից:
\end{hint}

% 17.10
\item Ապացուցեք. Ամեն մի կոմպակտ հաուսդորֆյան տարածություն ռեգուլյար տարածություն է:

% 17.11
\item Ապացուցեք. Հաուսդորֆյան տարածության ցանկացած երկու չհատվող կոմպակտ ենթաբազմությունների համար գոյություն ունեն դրանց չհատվող շրջակայքեր:

% 17.12
\item Ապացուցեք. Ամեն մի կոմպակտ հաուսդորֆյան տարածություն նորմալ տարածություն է:
    
\end{enumerate}
